%@ pages 172-184
%@ status translated
% Original section 17, "Préfaisceaux de Quillen", pp. 172--184.

\clearpage
\section{Quillen presheaves}
\label{quipage}

It was in studying the example of the stack of modules of complexes (cf.\ \S21
below) that we felt the need to consider ``presheaves of cmc''.
A point of view which is currently spreading is that the ``good'' morphisms
between closed model categories are the {\em left Quillen functors} (a notion
recalled in \S4)---cf.\ for example the book of Dwyer--Hirschhorn--Kan
\cite{DHK}, or also \cite{HoveyBook}, where Hovey defines the {\em $2$-category
of closed model categories}, whose objects are the closed model categories and
whose morphisms are the left Quillen functors.
One is thus naturally led to the following notion.


Let $\Xx$ be a category. A {\em left Quillen presheaf} (resp.\ a {\em right
Quillen presheaf}) on $\Xx$ is a presheaf of categories ${\bf M}$, in which each
${\bf M}_X$ is equipped with a structure of closed model category, such that the
restriction functors are left (resp.\ right) Quillen functors.
\footnote{
It would perhaps be preferable to speak of a {\em Quillen $1$-stack}, since
the values are $1$-categories and one could envisage a coherence problem. In
fact, since the notion of closed model category is stable under equivalence of
the underlying category, one can always strictify and speak only of presheaves
of categories. There could also be a coherence problem for the choice of
adjoints of the restriction functors. We ignore these problems, considering
that if they exist they are of neither geometric nor topological interest
(just as we ignore the subtleties of cardinality questions).}
Note that the restriction functors do not necessarily preserve all the
structure (they preserve the cofibrations and trivial cofibrations, but not
necessarily the fibrations nor the weak equivalences which are not
cofibrations).

If one draws a connection between the notions of topos and of cmc, the notions
of left or right Quillen presheaf are analogous to the notions of ``category
fibred in toposes'' and ``category fibred in duals of toposes'' of SGA 4
(\cite{SGA4b} Exposé Vbis, Def.~1.2.1, 1.2.2).

If ${\bf M}$ is a left Quillen presheaf, let ${\bf M}^{\ast}$ be the presheaf
which to each $x\in \Xx$ associates the closed model category ${\bf
M}^{\ast}(x)$ dual to ${\bf M} (x)$. Then ${\bf M}^{\ast}$ is a right Quillen
presheaf. This allows us to take an interest only in left Quillen presheaves. A
little more interesting: let ${\bf M}^{\dag}$ be the presheaf on the opposite
category $\Xx ^o$, with the same closed model categories ${\bf M}(x)$ as
values, but obtained by using the right adjoints as restriction morphisms.
Then ${\bf M}^{\dag}$ is also a (right) Quillen presheaf.

Let ${\bf M}$ be a Quillen presheaf over $\Xx$. The cofibrant objects define a
subpresheaf of categories ${\bf M}_c$ and its subpresheaf $W{\bf M}_c$ in which
the morphisms are the weak equivalences between cofibrant objects (which are
indeed preserved by the restrictions). One defines a presheaf of simplicial
categories $L({\bf M})$ over $\Xx$ by taking for
$
L({\bf M})(X)
$
the Dwyer--Kan localised simplicial category
$$
L(\mm )(X):= L({\bf M}_c(X), W{\bf M}_c(X)).
$$
As usual, one regards this presheaf as a Segal $1$-prestack on $\Xx$.

To simplify the later notation, it is more convenient to take the ``covariant''
point of view. From now on we consider a category $Y$ and a left Quillen
presheaf $\mm$ on $Y^o$; it is thus a family of cmc's $\mm (y)$ with left
Quillen restriction functors, which are covariant.



Let $\mm$ be a left Quillen presheaf on $Y^o$; if $f: y\rightarrow z$ is a
morphism of $Y$, let us denote by
$$
res_f: \mm (y)\rightarrow \mm (z)
$$
the left Quillen restriction functor; and let $res_f^{\ast}$ denote its right
adjoint. One has the ``fibred category'' (cf.\ the preceding section)
$$
\int _Y\mm \rightarrow Y
$$
and the category of sections
$$
Sect (Y,\int _Y\mm ).
$$
An object $\sigma$ of $Sect (Y,\int _Y\mm )$ consists of a collection of
objects $\sigma (y)\in
\mm (y)$ for $y\in Y$, together with morphisms
$$
\sigma (f): res_f\sigma (y)\rightarrow \sigma (z),
$$
which one may interpret by adjunction:
$$
\sigma (f): \sigma (y)\rightarrow res_f^{\ast}\sigma (z).
$$
These morphisms are subject to an associativity constraint.
A morphism $a:\sigma \rightarrow \tau$ in $Sect (Y,\int _Y\mm )$
is a collection of morphisms
$$
a(y): \sigma (y)\rightarrow \tau (y)\;\;\;\; \mbox{in}\;\;\; \mm (y)
$$
subject to the usual naturality condition.
We shall say that a morphism $a$ is a {\em weak equivalence} if, for each $y$,
$a(y)$ is a weak equivalence in $\mm (y)$.

One may observe that if $\mm = \underline{M}$ is the constant presheaf taking as
value the fixed closed model category $M$ (which is indeed a left Quillen
presheaf), then one has
$$
\int _Y\underline{M} = Y\times M
$$
and
$$
Sect (Y,\int _Y\underline{M}) = M^Y.
$$
Recall that one has at one's disposal several structures of closed model
category on $M^Y$ under various hypotheses concerning $M$ and $Y$.
\newline
(I)\,\, First there is the structure generated by Hirschhorn's cofibrations, of
Bousfield--Kan type (we had called it of ``type HBKQ'' in \S5);
this one exists as soon as $M$ is
generated by cofibrations (Hirschhorn \cite{Hirschhorn}). The fibrations are
the object-by-object fibrations.
\newline
(II)\,\, Next there is the structure which Dwyer--Hirschhorn--Kan qualify as
``unusual'' but which is the structure usually used in $K$-theory, and
apparently due to K.~Brown, Joyal and Jardine in the case of diagrams of
simplicial sets (Hirschhorn calls it the ``Heller structure''); here the
cofibrations are the object-by-object cofibrations. The existence of this
structure seems to coincide roughly with that of the structure (I) of type
HBKQ, but we have not found in the literature a general criterion for this
existence (cf.\ however the book of Goerss--Jardine \cite{JardineGoerssBook}).
\newline
(III)\,\, Finally, if $Y$ is a Reedy category, there is a model structure on
$M^Y$ with no further hypothesis on $M$ (at most, if the degree function is
indexed by an ordinal larger than $\omega$, one must assume that limits of that
size exist in $M$). This structure is in general different from the other two
(for the definition, see the proof of the theorem below). For this {\em Reedy
structure} see \cite{Reedy}, \cite{Hirschhorn}, \cite{DHK},
\cite{JardineGoerssBook}. The essential part of the notion of Reedy fibration
also appears in the notion of hypercovering expounded by Verdier in SGA 4
(\cite{SGA4b}, Exposé V, \S7, definition 7.3.1.1 (H3)).

If $\mm$ is a left Quillen presheaf on $Y^o$, one has the analogous structures
on $Sect (Y,\int _Y\mm )$:

\begin{theorem}
\label{sectionscmf}
Let $\mm$ be a left Quillen presheaf on $Y^o$. Then one has the following
structures of closed model category on $Sect (Y,\int
_Y\mm )$ (the weak equivalences are always the morphisms which are
object-by-object weak equivalences):
\newline
(I)\,\, The structure of type HBKQ in which the fibrations are the morphisms
$a:\sigma \rightarrow \tau $ such that $a(y)$ is a fibration for every $y\in
Y$: it exists if each $\mm (y)$ is generated by cofibrations;
\newline
(II)\,\, The structure of type Brown--Jardine--Heller in which the cofibrations
are the morphisms $a:\sigma \rightarrow \tau$ such that $a(y)$ is a cofibration
for every $y\in Y$,
%%%%%%%%%%%%%%%%%%%%%%%%%%%%%%%%%%%%%%%%%%%%%
%%
%% old condition: unusable because the retractions will never be functorial ....
%%
%% if each $\mm (y)$ is generated by cofibrations with
%% generating sets $I(y)$, $J(y)$ such that $res _f(I(y))\subset I(z)$
%% (resp.\ $res _f(J(y))\subset J(z)$) for every arrow $f: y\rightarrow z$,
%% and if
%% the $I(y)$-cofibrations are functorially (in $y$) retracts of
%% regular $I(y)$-cofibrations; and
%%%%%%%%%%%%%%%%%%%%%%%%%%%%%%%%%%%%%%%%%%%%%
%% here is the new condition:
%%%%%%%%%
provided that this class of cofibrations (resp.\ the cofibrations $\cap$ the
weak equivalences) admits a generating set $I$ (resp.\ $J$) allowing the small
object argument
(see before \ref{dhklemme}),
and that moreover each $\mm (y)$ is generated by cofibrations; and
\newline
(III)\,\, The structure of Reedy type (cf.\ the definition
below), if $Y$ is a Reedy category.
\end{theorem}
Compare with
Proposition 1.2.12 of SGA 4 \cite{SGA4b} Exposé Vbis.

\noindent
{\em Proof:}
It may be useful to review the proof of the statement \ref{dhklemme}
above, but for (I) we advise the reader to turn directly to the references
\cite{Hirschhorn} \cite{DHK} instead of keeping to \ref{dhklemme}, since we
shall use here directly arguments of the type of \cite{DHK}.

Note first that the set $W$ of weak equivalences automatically satisfies
closure under retracts and the ``three for the price of two'' property.

In (I) and (II) the condition that each $\mm (y)$ be generated by
cofibrations (and
hence admit small limits and colimits) implies that
$Sect (Y,\int
_Y\mm )$ admits small limits and colimits.


For (I) the proof is the same as in \cite{Hirschhorn},
\cite{DHK}.
Let $\sigma$ be a section and $i: \sigma (y)\rightarrow u$ a cofibration in
$\mm (y)$. One then obtains the {\em cofibration freely generated by $i$},
$Lib(i): \sigma \rightarrow \nu$, defined by
$$
\nu (z) = \coprod _{f: y\rightarrow z} res_f(u) \sqcup ^{\coprod_{f}
res_f(\sigma (y))}\sigma (z).
$$
Let $I$ (resp.\ $J$) be the set of cofibrations
of the form $Lib(i)$ with $i\in I(y)$, $y\in Y$ (resp.\ $i\in J(y)$, $y\in
Y$).
The $I$-injectives (resp.\
$J$-injectives) are just the morphisms which are object-by-object fibrations
(resp.\ trivial fibrations); and one checks easily that the recognition
criterion (``Recognition lemma'' \cite{Hirschhorn} 13.3.1, \cite{DHK} 8.1)
is satisfied. \footnote{Essentially the
only thing to check---as is indicated in
(\cite{DHK}, proof of 48.7: the model category structure on
$S-Cat$, p.~61), see also our Lemma \ref{dhklemme}---is that the colimit of
regular $J$-cofibrations is again in $W$. In the present case
this is
true object-by-object over each object of $Y$.}

%% For (II) one takes as generating sets $I$ the set of
%% morphisms $a$ such that $a(y)\in I(y)$ for every $y\in Y$ (resp.\ $J$ the
%% set
%% of morphisms $a$ such that $a(y)\in J(y)$). Since these sets
%% contain the
%% sets of (I) above, one obtains that the $I$-injectives are in $W$.
%% On the other hand the hypothesis implies that a morphism which is
%% object-by-object an
%% $I(y)$-cofibration for every $y$ is a retract of a regular $I$-cofibration,
%% hence the cofibrations are the morphisms which are cofibrations
%% object-by-object. The coproduct or sequential colimit of morphisms in
%% $J$ is
%% a weak equivalence, which leads to the verification of the recognition
%% criterion (\cite{Hirschhorn} 13.3.1, \cite{DHK} 8.1), and gives (II).

For (II) we rely on Lemma \ref{dhklemme}. Recall that the condition that each
$\mm (y)$
is generated by cofibrations gives us the properties of
\ref{dhklemme} (0)--(7) for each
$\mm (y)$. As a consequence one obtains immediately the properties (0), (1),
(2),
(6) and (7) for $Sect (Y,\int
_Y\mm )$. One may remark that the cofibrations of the structure (I)
above are
cofibrations here as well, so the lifting property for these
cofibrations implies that a
morphism is object-by-object a weak equivalence. This gives (3).
Finally (4) and (5) form
the first part of the hypothesis here, so by Lemma \ref{dhklemme}
one obtains (II).

{\em Remark:} Just as for the application of Lemma \ref{dhklemme}, one
may say that the part
of the hypothesis (II) which corresponds to \ref{dhklemme} (4) and (5) would be
a consequence of the
set-theoretically reasonable character of the notions of cofibration and
weak equivalence,
in all the examples one will encounter. See Jardine \cite{Jardine} for the
technique needed to
give a rigorous proof of this type of condition. In view of this
remark, one may
consider that the hypothesis of (II) reduces to asking (as for
(I)) simply that the
$\mm (y)$ be generated by cofibrations.

\subnumero{Reedy categories}


Before tackling part (III) of Theorem \ref{sectionscmf},
we recall a few facts about ``Reedy categories''. The
references are
\cite{Reedy}, \cite{BousfieldKan}, \cite{Hirschhorn}, \cite{DHK},
\cite{JardineGoerssBook}, see also \cite{DwyerKan3}.

The interplay between the ``latch'' and the ``match'' which follows is only a
systematic way, with an abstract notation, of taking into account the fact that
there are degeneracy morphisms in $\Delta$ which go in the opposite direction
from the ``face'' morphisms. In all applications one can probably consider,
instead of simplicial objects, objects parametrised by the subcategory
$\direct{\Delta}$ generated by the face morphisms. In that case one would only
have to consider the ``Latch'' part of the discussion below and many arguments
would become easier.


We first recall the definition: a {\em Reedy category}
is a category $Y$ equipped with two subcategories (with the same objects
as $Y$) called respectively the {\em direct subcategory} and the {\em
inverse subcategory}
$$
\direct{Y}\subset Y \;\;\; \mbox{and} \;\;\; \inverse{Y}\subset Y,
$$
such that there exists a {\em degree} function (to an ordinal which will most
often be $\omega$) such that the morphisms of $\direct{Y}$ other than the
identities strictly increase the degree, the morphisms of $\inverse{Y}$ other
than the identities strictly decrease the degree, and every morphism $f$
factors uniquely
$$
f= \direct{f}\inverse{f} \;\;\; \mbox{with} \;\;\; \direct{f}\in \direct{Y},
\;\; \inverse{f}\in \inverse{Y}.
$$
We specify here that we write compositions in the usual direction, so in this
formula the source of $\direct{f}$ (which is the target of $\inverse{f}$) has
degree $\leq$ the minimum of the degrees of the source and of the target of
$f$; and in case of equality $f$ is either direct or inverse.

The classical examples are the category $\Delta$ and its opposite $\Delta
^o$, \cite{BousfieldKan}, \cite{Reedy}.
The example we shall use is the category $\Delta K$ of simplices of a
simplicial set $K$ (or its opposite $\Delta ^oK$). This example appears
in \cite{Hirschhorn}, \cite{DHK}, etc. More particularly we shall use
the case of $\Delta \nu \Yy$ where $\nu \Zz$ is the nerve of a
$1$-category $\Zz$. It is the category denoted $\beta (\Zz )$ in \S16
above.




For an object $y\in Y$ one defines the {\em latching} ({\em relâchantes})
categories
\footnote{
The word ``relâchant'' is not a strict translation of
``latching'' but has the same ring to it and seems
to correspond well to the situation.}
$Latch(y)$
and the {\em matching} ({\em appariantes}) categories $Match(y)$.
In fact one defines
$$
Latch(y)+\{ y\} := \direct{Y}/y
$$
(the objects of $\direct{Y}$ over $y$)
and one sets
$$
Latch(y):= \direct{Y}/y -\{ y\} \subset Latch(y)+\{ y\} .
$$
Also
$$
\{ y\} + Match(y):= y/\inverse{Y}
$$
(the objects of $\inverse{Y}$ under $y$)
and
$$
Match(y):= y/\inverse{Y} -\{ y\} \subset \{ y\} + Match(y).
$$
One denotes by
$$
Latch(y)+ \{ y\} + Match(y)
$$
the category obtained by attaching $Latch(y)+\{ y\} $ to
$\{ y\} + Match(y) $ along the object $y$, with exactly one morphism between
the objects of $Latch(y)$ and $Match(y)$. One denotes by
$$
Latch(y)+Match(y)\subset Latch(y)+ \{ y\} + Match(y)
$$
the subcategory obtained by removing $y$.
The morphism of Segal $1$-precategories
$$
(Latch(y)+\{ y\} ) \cup ^{\{ y\} } (\{ y\} + Match(y))\rightarrow
Latch(y)+ \{ y\} + Match(y)
$$
is a weak equivalence (we leave the proof, which is similar to that
of \ref{reedydecomp} below, to the reader). In particular, for a
fibrant Segal $n$-category $C'$ a morphism
$$
Latch(y)+ \{ y\} + Match(y)\rightarrow C'
$$
is essentially the same thing as two morphisms
$$
Latch(y)+\{ y\} \rightarrow C', \;\;\; \{ y\} + Match(y) \rightarrow C'
$$
which send $y$ to the same object of $C'$.

For a fixed degree $k$, one denotes by $F^kY\subset Y$ the subcategory
of objects of degree $\leq k$. If $y$ has degree $k$ then one has a
morphism
$$
Latch(y)+Match(y) \rightarrow F^{k-1}Y.
$$
One has that $Y$ is the filtered colimit of the $F^kY$.

The description of diagrams on a Reedy category of \cite{DHK},
\cite{JardineGoerssBook}, \cite{BousfieldKan},
\cite{DwyerE2}, \cite{BousfieldFriedlander}, \cite{Reedy}, \cite{DwyerKan3},
and particularly Hirschhorn (\cite{Hirschhorn} Theorem 16.2.12) where this
appears very clearly, amounts essentially to the following lemma.

\begin{lemma}
\label{reedydecomp}
For every $d$, the morphism of Segal $1$-precategories
$$
F^{d-1}Y \cup ^{\coprod _{deg(y)=d}Latch(y)+Match(y)}\coprod_{deg(y)=d}
Latch(y)+ \{ y\} + Match(y)
$$
$$
\rightarrow F^dY
$$
is a weak equivalence.
\end{lemma}

This lemma possibly generalises the above references by taking into account
(via the
notion of Segal $1$-precategory) the higher homotopies. We shall need
this aspect in the following chapters. For part
(III) of Theorem \ref{sectionscmf}, the version
(\cite{Hirschhorn} Theorem 16.2.12)---which is the statement obtained from our
lemma by
applying the truncation $\tau_{\leq 1}$---would suffice. In view of this, we
postpone the
proof until the end of the chapter.



Lemma \ref{reedydecomp} says that for a fibrant Segal $n$-category $C'$, to
give a morphism
$f_k:F^kY\rightarrow C'$ amounts to giving a morphism
$f_{k-1}: F^{k-1}Y\rightarrow C'$ and, for every object $y$, extensions of the
morphisms
$$
f_{k-1}|_{Latch(y)} \;\; \mbox{into}\;\; Latch(y)+\{ y\} \rightarrow C',
$$
and
$$
f_{k-1}|_{Match(y)} \;\; \mbox{into}\;\; Match(y)+\{ y\} \rightarrow C',
$$
taking the same value on $y$,
and such that the induced morphism (essentially well defined)
$$
Latch(y)+ \{ y\} + Match(y)\rightarrow C'
$$
has $f_{k-1}$ as its restriction to $Latch(y)+Match(y)$.
This interpretation comes from Hirschhorn \cite{Hirschhorn}.


\subnumero{The structure (III) of \texorpdfstring{Theorem~\ref{sectionscmf}}{Theorem}}

One must say what the fibrations and cofibrations are.
Let $\sigma$ be a section and $y\in Y$.
Let $latch(\sigma , y)$ and $match(\sigma , y)$ in $\mm (y)$ be the {\em
latching} and {\em matching objects} of $\sigma$. These objects are defined
respectively
as the colimit over $Latch(y)$ of the functor
$$
res_{Latch(y)\rightarrow y} \circ \sigma : Latch(y)\rightarrow \mm (y)
$$
where
$$
res_{Latch(y)\rightarrow y} : \int _{Latch(y)}\mm |_{Latch(y)} \rightarrow \mm
(y)
$$
is the morphism given by the restrictions $res_f$ for $f: z\rightarrow y$;
and the limit over $Match(y)$ of the functor
$$
res_{y\rightarrow Match(y)} ^{\ast}\circ \sigma : Match(y)\rightarrow \mm
(y)
$$
where
$$
res_{y\rightarrow Match(y)}^{\ast} : \int _{Match(y)}\mm |_{Match(y)}
\rightarrow \mm (y)
$$
is the morphism given by the adjoints of the restrictions
$res_g^{\ast}$ for $g: y\rightarrow w$.
To define $res_{y\rightarrow Match(y)}^{\ast}$ one notes that one has, for
$$
y\stackrel{g}{\rightarrow} w \stackrel{h}{\rightarrow} x,
$$
a natural transformation
$$
res_g^{\ast} \rightarrow res _{hg}^{\ast} \circ res_h,
$$
which comes from the adjunction morphism
$$
1\rightarrow res_h^{\ast} \circ res_h.
$$
One has a factorisation
$$
latch(\sigma ,y)\rightarrow \sigma (y) \rightarrow match(\sigma , y).
$$
We shall say (following \cite{Reedy} \cite{Hirschhorn} \cite{DHK} {\em et al.})
that
$\sigma$ is {\em (Reedy) cofibrant} if the morphisms $latch(\sigma
,y)\rightarrow \sigma (y)$ are cofibrations in $\mm (y)$, and that
$\sigma$
is {\em (Reedy) fibrant} if
the morphisms $\sigma (y)\rightarrow match(\sigma ,y)$ are fibrations in
$\mm (y)$. More generally for $a:\sigma \rightarrow \tau$
one obtains the relative latching object
$$
latch(a,y):= latch(\tau , y) \cup ^{latch(\sigma , y)}\sigma (y)
$$
with morphism
$$
latch (a,y)\rightarrow \tau (y);
$$
we shall say that $a$ is a {\em (Reedy) cofibration} if this morphism is a
cofibration in $\mm (y)$. Dually one obtains the relative matching object
$$
match (a,y):= match (\sigma ,y)\times _{match(\tau , y)}\tau (y)
$$
with morphism
$$
\sigma (y)\rightarrow match(a,y);
$$
and we shall say that $a$ is a {\em (Reedy) fibration} if this morphism is a
fibration in $\mm (y)$.

The proof that these three classes of morphisms form a structure of
closed model
category is the same as in the case of the constant presheaf
$\mm =\underline{M}$, cf.\ \cite{Reedy} \cite{Hirschhorn} {\em et al.};
we leave this verification to the reader.

This completes the proof of (III) and hence of Theorem
\ref{sectionscmf}.
\eop

\subnumero{Proof of \texorpdfstring{Lemma~\ref{reedydecomp}}{Lemma}}

We keep the notation of the lemma to be proved.

One may assume that $Y=F^dY$ and one sets $Z:= F^{d-1}Y$. Let $A^0$ denote the
coproduct of the statement. It is a simplicial set
(subset of $\nu Y$), which one regards as a Segal $1$-precategory. We shall
obtain $\nu Y$ as the limit of a sequence (possibly transfinite if the number
of objects of $Y$ is more than countable)
$$
\ldots A^i\subset A^{i+1}\ldots
$$
of simplicial subsets of $\nu Y$. Each time one adds a
``regular simplex'' (definition below), and one proves that the
$A^i\subset A^{i+1}$ are trivial cofibrations of Segal $1$-precategories.

Recall that the elements of $\nu (Y)_p$ are the composable sequences
$$
(y,f) =(y_0,\ldots , y_p; f_1,\ldots , f_p)
$$
with $f_i:
y_{i-1}\rightarrow y_i$ in $Y$. We shall say that such an element is a {\em
regular simplex} if none of the $f_i$ is the identity, and if for every
$1\leq i \leq p$, either $f_i \in \direct{Y}$, or $f_i \in \inverse{Y}$.
For a regular simplex, we shall call a {\em peak} any vertex, i.e.\ object
$y_i$ with $f_{i-1}\in \direct{Y}$ and $f_i\in \inverse{Y}$, which is
equivalent
to the condition $deg(y_{i-1}) < deg(y_i) > deg(y_{i+1})$. We shall call a {\em
valley} any object $y_i$ such that $f_{i-1}\in \inverse{Y}$ and $f_i\in
\direct{Y}$, which is equivalent to the condition $deg(y_{i-1}) > deg(y_i) <
deg(y_{i+1})$. The Reedy condition implies that if $y_i$ is a valley,
then the simplex in question is the only regular simplex of length $\leq
p$ containing the face obtained by removing $y_i$.

For an element $(y,f)\in \nu (Y)_p$ one obtains a morphism
$(y,f):h(p)\rightarrow Y$ where $h(p)$ is the simplicial set represented
by $p$, i.e.\ the standard $p$-simplex.
We shall call the ``image of a simplex $(y,f)$'' the image
of this morphism.
The image contains
all the ``hyperface'' simplices of $(y,f)$, which are all the simplices
obtained by removing some of the $y_i$, as well as all their
degeneracies.

One chooses a total order $\prec$ on the set of regular simplices of
$\nu Y$ not contained in $A^0$, subject to the following conditions: if
$$
(y_0,\ldots , y_p; f_1,\ldots , f_p) \prec
(y'_0,\ldots , y'_q; f'_1,\ldots , f'_q)
$$
then
$$
p\leq q
$$
and, if $p=q$ then
$$
\sum _{i=0}^p deg (y_i) \leq \sum _{i=0}^p deg(y'_i).
$$
Now one starts from $A^0$ and adds the images of $h(p)$ for the
regular simplices $(y,f): h(p)\rightarrow \nu (Y)$, one by one, following
the total order $\prec$ (at limit ordinals one inserts an object which is the
colimit of what comes before, then one starts again by adding the next
simplex). One obtains a sequence $A^i\subset \nu (Y)$. We shall show that a
cofibration $A^i\rightarrow A^{i+1}$ which corresponds to the addition of a
regular simplex $(y_0,\ldots , y_p; f_1,\ldots , f_p)$ is a
weak equivalence of Segal $1$-precategories. This will give the lemma
(there is a filtered colimit argument to be made for limit ordinals and
at the end of the process, which we leave to the reader).


We shall prove, by induction on $i$, the following statement:

$(\ast )$ a degenerate simplex
is in $A^i$ if and only if the corresponding non-degenerate simplex
is in $A^i$.

Let $A^i\rightarrow A^{i+1}$ be a cofibration which corresponds to
the addition of a regular simplex $(y_0,\ldots , y_p; f_1,\ldots , f_p)$, and
let $h(p)\rightarrow \nu (Y)$ be the morphism corresponding to this simplex.
If one sets
$$
U:= h(p)\times _{A^i} \nu (Y) \subset h(p),
$$
then one has
$$
A^{i+1} = h(p)\cup ^U A^i.
$$
It is therefore a matter of proving that $U\rightarrow h(p)$ is a trivial
cofibration of Segal $1$-precategories. Note first that by hypothesis $(\ast )$ for
$A^i$, a degenerate element of $h(p)_q$ is in $U$ if and only
if the corresponding non-degenerate element is in $U$. It is therefore
a matter of understanding which hyperfaces of $h(p)$ are in $U$. One has the
following characterisation:

$(\ast \ast )$ \, {\em A hyperface of $h(p)$ which corresponds to a
subset $J\subset \{ 0,\ldots , p\}$ belongs to $U$ if and only
if there exists $i\in \{ 0,\ldots , p\}$ with $i\not \in J$ and such that $y_i$
is not a
valley of $(y_0,\ldots , y_p; f_1,\ldots , f_p)$.}


Let us prove this. One considers the hyperface corresponding to
$$
J= \{j_0,\ldots ,
j_q\} \subset \{ 0,\ldots , p\}
$$
(strict inclusion). This corresponds to
the element
$$
J^{\ast}(y,f):=(y_{j_0}, \ldots , y_{j_q}, f'_1,\ldots , f'_q)\in \nu Y
$$
where
the $f'_j$ are composites of the $f_i$. Let $(z,g)=(z_0,\ldots , z_r;
g_1,\ldots , g_r)$ be the regular simplex obtained by replacing in
$J^{\ast}(y,f)$ every arrow $f'_k: y_{j_{k-1}}\rightarrow y_{j_k}$ which is
not
in $\direct{Y}$ nor in $\inverse{Y}$, by the pair of arrows
$$
y_{j_{k-1}}\stackrel{g'}{\rightarrow}z\stackrel{g''}{\rightarrow} y_{j_k}
$$
with $g'\in \inverse{Y}$ and $g''\in \direct{Y}$.
One leaves unchanged the
arrows $f'_k$ which are either in $\direct{Y}$ or in
$\inverse{Y}$.
Also one contracts all the identities $f'_k$ which appear in
$J^{\ast}(y,f)$. One notes that $J^{\ast}(y,f)$ is in the image of $(z,g)$.

An arrow $f'_k$ which is replaced (or contracted) in the
above procedure necessarily comes from a composite of at least two arrows
$f_i$. Hence the length $r$ of $(z,g)$ is $\leq p$. On the other hand, if $r=p$
then
$$
\sum _{i=0}^p deg (z_i) \leq \sum _{i=0}^p deg (y_i).
$$
Indeed, if $f'_k=f_{i+1}f_i = g_{i+1}g_i$ then $y_i$ is either a peak
or a valley: in the contrary case, $f'_k$ is either in
$\direct{Y}$
or in $\inverse{Y}$ and is not replaced in $(z,g)$ (and $f'_k$
is not an identity either, so it is not contracted either).
If $y_i$ is a valley then $z_i=y_i$; if $y_i$ is a peak, then in this case
$z_i$ is a valley of $(z,g)$ and $deg(z_i) < deg (y_i)$.

By the preceding paragraph, if the complement of $J$ in $\{
0,\ldots , p\}$ contains a peak $y_i$ then one of the two
inequalities is strict and (by what we imposed on the order
$\prec$) one has $(z,g)\prec (y,f)$ and hence $J^{\ast}(y,f)\in A^i$. If the
complement
of $J$ contains a $y_i$ which is neither a peak nor a valley, then since one of
the arrows $f'_k$ which is a composite of at least two arrows $f_i$ is
not replaced by two arrows $g$, one has $r<p$ and again
$J^{\ast}(y,f)\in A^i$.
This proves one direction of the statement $(\ast \ast )$
(one should note that, in the above argument, one may land, without saying so,
on a
simplex which is in $A^0$ instead of being added in the sequence organised
by $\prec$).


Moreover, if there is a contraction of some $f'_k= 1$ then the length
decreases: $ r<p$ and $J^{\ast}(y,f)\in A^i$. One sees that every new
simplex
of the form $J^{\ast}(y,f)$ in $A^{i+1}$ which is degenerate belongs
already to $A^{i}$; and the other degenerate simplices
come from the degeneracies of simplices of the form
$J^{\ast}(y,f)$. This proves $(\ast )$ for $A^{i+1}$.

To complete the proof of $(\ast \ast )$, suppose that the complement
of $J$ consists entirely of valleys of $(y,f)$. In this case, for every
$i\not \in J$, one has $i-1\in J$ and $i+1\in J$, and one of the morphisms in
$J^{\ast}(y,f)$ is
$$
f'_k=f_{i+1}f_i: y_{i-1}\rightarrow y_{i+1}.
$$
One notes that $f'_k$ is not in $\direct{Y}$ nor in $\inverse{Y}$.
Hence, if the simplex $J^{\ast}(y,f)$ comes from the image of a regular
simplex
$(w,h)$ then the arrow $f'_k$ is necessarily decomposed in $(w,h)$
as a product of at least two arrows $h$. It follows that the length of
$(w,h)$
is at least $p$. In case of equality, each of the arrows $f'_k$ will be
decomposed into a product of the form $h_{i+1} h_i$. Here there are two
options: either $w_i$ is a valley, in which case one has $w_i=y_i$ by
the uniqueness of
the decomposition in the axiom for Reedy categories; or $w_i$ is
a peak, in which case one has $deg (w_i)> deg(y_i)$. Hence, finally, there are
three cases:
either $(w,h)$ has length $>p$; or $(w,h)$ has length $p$ but
$$
\sum _{i=0}^p deg(w_i)>  \sum _{i=0}^p deg(y_i);
$$
or $(w,h)=(y,f)$. In all cases $(w,h)$ does not precede $(y,f)$ for
the order $\prec$ and hence $J^{\ast}(y,f)$ does not belong to $A^i$.

In fact, for this last sentence one must also justify that
$J^{\ast}(y,f)$ does not belong to $A^0$. One notes that $(y,f)$ has
at least one peak of degree $d$ (i.e.\ in $Y$ but not in $Z$), and this peak
persists in
$J^{\ast}(y,f)$. The only simplices of this form in $A^0$ are
the regular simplices with a single peak (i.e.\ the simplices added in
the coproduct of the statement of the lemma); but by removing valleys one has
made
$J^{\ast}(y,f)$ irregular and one has therefore $J^{\ast}(y,f)\not \in A^0$.
This
proves $(\ast \ast )$.

The statement $(\ast \ast )$ says that $U$ is the union of all the
faces of $h(p)$ obtained by removing vertices other than valleys.
This implies
that the morphism
$$
U\rightarrow
h(p)
$$
is a weak equivalence of Segal $1$-precategories, by the following
lemma (which we isolate because the statement could be of interest in its own
right). This will complete the
proof of
Lemma \ref{reedydecomp}.
\eop

\begin{lemma}
\label{faces}
Let $h(p)$ denote the simplicial set ``standard $p$-simplex'', and let
$f:U\hookrightarrow h(p)$ be an inclusion such that $U$ is the union of a
certain
number of faces of $h(p)$. Suppose that $U$ contains the first and the
last face, and does not contain all the others; then the morphism
$f$ is a weak equivalence of Segal $1$-precategories, i.e.\ the
induced morphism
$$
SeCat(U)\rightarrow SeCat(h(p))\cong I^{(p)}
$$
is an equivalence.
\end{lemma}
{\em Proof:}
One argues by induction on $p$; suppose therefore
the statement known for $p-1$.
Let us write
$$
U= V^0\cup V^p \cup V^{b_1}\cup \ldots \cup V^{b_k}
$$
as a union (not disjoint!) of faces, where $V^j$ is the face (simplicial
set isomorphic to $h(p-1)$) obtained by removing the $j^{\rm th}$
vertex. One fixes
$a\in \{ 1,\ldots , p-1\}$ with $b_i\neq a$, i.e.\ the face $V^a$
does not appear
in $U$. For $0\leq j\leq k$ set
$$
U^j:= V^0\cup V^p \cup V^{b_1}\cup \ldots \cup V^{b_j}\subset U.
$$
In particular $U^0= V^0\cup V^p$, and $U^k=U$. Note first that
$U^0\rightarrow h(p)$ is an equivalence (from here to the end of the proof,
this will mean that it is a weak equivalence of Segal $1$-precategories).
To see this, let $W(p)\subset h(p)$ be the union of the principal
edges,
i.e.\ of the edges of the form $\{ i-1, i\}$. The morphism
$W(p)\rightarrow h(p)$ is an equivalence (indeed, $h(p)$ is obtained
from $W$ by one application of the operation $Gen[p]$, see \S2;
see also the discussion in \cite{limits} \S2.4.8).
Now $V^0\cap W(p)$ is of the form $W(p-1)\subset V^0=h(p-1)$.
Hence one may express $Q(p):=V^0\cup W(p)$ as a coproduct:
$$
Q(p):=V^0\cup W(p) = V^0 \cup ^{W(p-1)}W(p)
$$
where the first morphism is a trivial cofibration (of Segal
$1$-precategories). Hence the morphism $W(p)\rightarrow Q(p)$ is a trivial
cofibration and one obtains that $Q(p)\rightarrow h(p)$ is an
equivalence (and it is a cofibration, hence it is a trivial cofibration).
Note that $Q(p)$ is the union of the first face with the last edge.
Now $Q(p) \cap V^p = Q(p-1)\subset V^p=h(p-1)$. Hence (by the result
just proved but for $Q(p-1)$) the second morphism in the
coproduct
$$
U^0= V^0\cup V^p = Q(p)\cup V^p = Q(p) \cup ^{Q(p-1)}V^p
$$
is a trivial cofibration. It follows that $Q(p)\rightarrow U^0$ is a
trivial cofibration and $U^0\rightarrow h(p)$ is an
equivalence.

Let us now show by induction on $j$ that $U^j\rightarrow h(p)$ is
an equivalence. We have just proved it for $j=0$, so one may suppose
that $j\geq 1$ and that it is known for $j-1$. One has
$$
U^j = U^{j-1} \cup ^{(U^{j-1}\cap V^{b_j})}V^{b_j}.
$$
To prove that $U^j\rightarrow h(p)$ is an equivalence it suffices to prove
that
$$
(U^{j-1}\cap V^{b_j})\rightarrow V^{b_j}
$$
is an equivalence. Now $V^{b_j} = h(p-1)$, and one denotes by
$0,\ldots , \widehat{b_j}, \ldots , p$
the vertices of this simplex. The simplicial subset
$(U^{j-1}\cap V^{b_j})$ is a union of faces of this $h(p-1)$. This
union
of faces contains the first and the last face (since $V^0\cup V^p\subset
U^{j-1}$), and does not contain the face obtained by removing the vertex
$a$. Hence the lemma for $p-1$ applies and one obtains that
$(U^{j-1}\cap V^{b_j})\rightarrow V^{b_j}$ is an equivalence, as required.
\eop
